\section{Nested shells}\label{sec:shells}
\stagebox{\statusbuilt}{\src{core/src/shells.rs}}

A bud is layered like an onion. Each petal layer lives on a thin egg-shaped
shell, and the bud's size follows from how many layers it packs and how thick
they are.

\subsection{Scaled radius}
Lengths are in units of the nominal core radius, with $+z$ pointing up. The
core itself, the central stamen and carpel zone, has radius $R_0 = s_c$, where
the core size $s_c$ is a species trait drawn per plant. Every shell is a level
set of the \emph{scaled radius}
\[
  \rt(\x) = \frac{|\mathbf q|}{g(c)}, \qquad
  \mathbf q = (x,\; y,\; z/\alpha), \qquad c = q_z / |\mathbf q|,
\]
so all shells are scaled copies of one outline: a spheroid stretched by the
aspect $\alpha$ along $z$, reshaped by the factor
\[
  g(c) = 1 + \eta_e\,P_3(c) + \eta_p\,\max(c, 0)^8, \qquad P_3(c) = \tfrac12(5c^3 - 3c).
\]
The egg term $\eta_e$ broadens the lower body and narrows the top ($c = 1$ is the
top pole); the point term $\eta_p$ draws the top out into a pointed tip, as in a
rose bud. Both are species traits drawn per plant, which gives the
plain spheroid when both are 0. Since the core is a shell too, it takes the
same outline; beyond about $|\eta_e| = 0.12$ the core turns visibly
triangular, and peony stays within $0 \le \eta_e \le 0.08$. An $\eta_p$ of about
0.3 gives a clearly pointed, rose-like bud; peony draws $0 \le \eta_p \le 0.1$.
The derivative of the shape factor is
\[
  g'(c) = \tfrac12\eta_e\,(15c^2 - 3) + \begin{cases} 8\eta_p\,c^7 & c > 0,\\ 0 & c \le 0,\end{cases}
\]
and $c^8$ is computed as $(c^2)^2(c^2)^2$ so that every platform rounds it the
same way (\cref{sec:determinism}). At the centre itself, $|\mathbf q| <
10^{-12}$, the scaled radius is taken as $|\mathbf q|$. With
$\mathbf u = \mathbf q/|\mathbf q|$, the gradient used for shell normals and by
the packing energy is
\[
  \frac{\partial\rt}{\partial\mathbf q} = \frac{\mathbf u}{g} - \frac{g'(c)}{g^2}\,\bigl(\mathbf e_z - c\,\mathbf u\bigr),
\]
with the $z$ component divided by $\alpha$ to give $\nabla\rt(\x)$.

\subsection{Slots}
Each layer owns a slot of width
\[
  w = \tmax\,(1 + \varrho),
\]
the thickest a petal can be plus some initial room ($\varrho$ is \texttt{crumple\_room}). Slots stack outward from the
core, so layer $k = 0,\dots,L-1$ has its mid-surface at
\[
  a_k = R_0 + w\,(k + \tfrac12),
\]
and the inner side of the sepal envelope starts at $R_{\mathrm{env}} = R_0 + wL$.

\begin{figure}[ht]
\centering
\begin{tikzpicture}[scale=1.6]
  \def\al{1.25}   % aspect
  \def\pb{25.8}   % phi_b in degrees (0.45 rad)
  \def\pg{5.7}    % phi_g in degrees (0.10 rad)
  \fill[yellow!25] (0,0) ellipse (1 and \al);
  \node[font=\small, fill=yellow!25, inner sep=1pt] at (0,0.3) {core $R_0$};
  \foreach \k/\c in {0/red!50, 1/orange!50, 2/magenta!40} {
    \pgfmathsetmacro{\a}{1 + 0.3*(\k+0.5)}
    \draw[\c, line width=3pt, domain=\pb:180-\pg, samples=60, variable=\p]
      plot ({\a*sin(\p)}, {-\al*\a*cos(\p)});
    \draw[\c, line width=3pt, domain=\pb:180-\pg, samples=60, variable=\p]
      plot ({-\a*sin(\p)}, {-\al*\a*cos(\p)});
  }
  \draw[dashed, gray] (0,0) ellipse (1.9 and 1.9*\al);
  \node[font=\small, gray, anchor=west] at (1.25,2.15) {envelope $R_{\mathrm{env}}$};
  \draw[{Stealth}-{Stealth}] (1.0,0) -- (1.3,0);
  \node[font=\scriptsize, above] at (1.15,0) {$w$};
  \draw[gray] (0,-2.6) -- (0,2.6);
  \node[font=\scriptsize, gray, anchor=west] at (0,2.55) {$z$};
  \draw[gray] (0,0) -- ({1.9*sin(\pb)}, {-1.9*\al*cos(\pb)});
  \draw[gray] (0,-0.55) arc[start angle=-90, end angle=-90+\pb, radius=0.55];
  \node[font=\scriptsize] at (0.22,-0.75) {$\phi_b$};
\end{tikzpicture}
\caption{Three layers in their slots, seen in cross-section (schematic: a
plain spheroid with $\eta_e=\eta_p=0$ and $s_c=1$). Petals attach on
the ring at polar angle $\phi_b$ from the bottom pole and stop $\phi_g$ short of
the top pole.}
\label{fig:shells}
\end{figure}

\subsection{Shell coordinates}
A point on the shell of scaled radius $a$ at azimuth $\theta$ and polar angle
$\phi$ (measured from the bottom pole) is
\[
  P(a,\theta,\phi) = r\,\bigl(\sin\phi\cos\theta,\; \sin\phi\sin\theta,\; -\alpha\cos\phi\bigr),
  \qquad r = a\,g(-\cos\phi).
\]
Petals occupy polar angles from the attachment ring $\phi_b$ up to
$\pi - \phi_g$. A fractional height $f\in[0,1]$ maps to
\[
  \phi(f) = \phi_b + f\,(\pi - \phi_g - \phi_b).
\]

\begin{remark}
The shells fix the starting layout only. During development the envelope
radius becomes a variable that petal pressure pushes outward, and the layer
radii move out with it in proportion (\cref{sec:packing}).
\end{remark}
